\documentclass[12pt]{article}  % Tipo de archivo
\usepackage[utf8]{inputenc}    % Codificación utf-8
\usepackage[spanish]{babel}    % Idioma español
\usepackage{graphicx}          % Necesario para imágenes

\title{
    \textsc{\Large Universidad Nacional Autónoma de México\\[0.2cm]
        {\normalsize Facultad de Ciencias, 2027-I}\\[0.2cm]
        {\normalsize Modelado y Programación}}\\[0.6cm]
    \rule{\linewidth}{0.4pt}\\[0.4cm]
    {\textsc{\Large {\color{maincolor}Proyecto:} Chat Multiusuario}}\\[0.1cm]
    \rule{\linewidth}{0.4pt}\\[1em]
}
\author{
    \textsc{Alumno:} Valeria Becerril Acosta\\[0.2cm]
    \textsc{Número de cuenta:} 323128489\\[0.4cm]
}
\date{\today\\[2.6cm]}

% PAQUETES DE ESTILO
\usepackage{xcolor}           % Para definir y usar colores
\usepackage{listings}         % Para el estilo de código mystyle
\usepackage{fancyhdr}         % Para los encabezados y pies de página fancy
\usepackage{titlesec}         % Para configurar \titleformat
\usepackage{enumitem}         % Para configurar \setlist
\usepackage[most]{tcolorbox}  % Para los cuadros de color
\usepackage{lmodern}          % Para la familia de fuentes
\usepackage[T1]{fontenc}      % Para codificación de fuente

% CONFIGURACIONES Y COLORES
\definecolor{maincolor}{RGB}{0, 51, 102}
\definecolor{codegreen}{rgb}{0,0.6,0}
\definecolor{codegray}{rgb}{0.5,0.5,0.5}
\definecolor{codepurple}{rgb}{0.58,0,0.82}
\definecolor{foldercolor}{RGB}{255, 193, 7}
\definecolor{pdfcolor}{RGB}{220, 53, 69}
\definecolor{readmecolor}{RGB}{40, 167, 69}
\definecolor{txtcolor}{RGB}{0, 123, 255}

% COMANDOS
\newcommand{\newItem}[3]{
  \item {\color{maincolor}\textbf{#1}}
    \begin{itemize}
      \item {\fontseries{b}\selectfont Petición del cliente:} #2
      \item {\fontseries{b}\selectfont Respuesta del servidor:} #3
    \end{itemize}
}

\newcommand{\uml}[3]{
    \begin{center}
        \begin{tcolorbox}[colback=maincolor!10, colframe=maincolor,
            colupper=black, fontupper=\footnotesize\ttfamily,
            boxrule=0.2mm, arc=2pt,
            left=3pt, right=3pt, top=1pt, bottom=1pt, width=2.5in,
            title={\centering\fontseries{b}\selectfont\footnotesize #1},
            colbacktitle=maincolor!35,
            coltitle=black,
            toptitle=4pt, bottomtitle=4pt]

            \vspace{4pt}
            \begin{itemize}[label={}, itemsep=-5pt]
                #2
            \end{itemize}
            \vspace{6pt}

            {\color{maincolor}\hrule height 0.2mm}

            \vspace{-3pt}
            \begin{itemize}[label={}, itemsep=-5pt]
                #3
            \end{itemize}
            \vspace{-6pt}
        \end{tcolorbox}
    \end{center}
}

\newcommand{\public}[1]{
  \item[$+$] #1
}

\newcommand{\private}[1]{
  \item[$-$] #1
}

\newcommand{\noAccess}[1]{
  \item[] #1
}

% ESTILO DE CÓDIGO
\lstdefinestyle{mystyle}{
    backgroundcolor=\color{maincolor!10},
    commentstyle=\color{codegreen},
    keywordstyle=\color{magenta},
    numberstyle=\tiny\color{codegray},
    stringstyle=\color{codepurple},
    basicstyle=\ttfamily\footnotesize,
    breakatwhitespace=false,
    breaklines=true,
    captionpos=b,
    keepspaces=true,
    numbers=left,
    numbersep=5pt,
    showspaces=false,
    showstringspaces=false,
    showtabs=false,
    tabsize=2
}
\lstset{style=mystyle}

% CONFIGURACIÓN DE ENCABEZADO Y PIE DE PÁGINA
\pagestyle{fancy}
\fancyhead[LO,L]{UNAM}
\fancyhead[RO,R]{\textsc{Proyecto}}
\fancyfoot[LO,L]{Facultad De Ciencias, 2027-I}
\fancyfoot[CO,C]{\thepage}
\fancyfoot[RO,R]{Modelado y Programación}
\renewcommand{\headrulewidth}{0.4pt}
\renewcommand{\footrulewidth}{0.4pt}

% CAMBIAR COLOR DE SECCIONES
\titleformat{\section}
{\color{maincolor}\normalfont\Large\bfseries}
{\thesection}{1em}{}

\titleformat{\subsection}
{\color{maincolor}\normalfont\large\bfseries}
{\thesubsection}{1em}{}

\titleformat{\subsubsection}
{\color{maincolor}\normalfont\normalsize\bfseries}
{\thesubsubsection}{1em}{}

% CAMBIAR COLOR DE ÍTEMS EN LISTAS
\setlist[enumerate]{label=\textcolor{maincolor}{\arabic*.}, left=1.5em}
\setlist[itemize,1]{label=\textcolor{maincolor}{\textbullet}, left=1.5em}
\setlist[itemize,2]{label=\textcolor{maincolor}{\textendash}, left=1.5em}


% INICIO DEL DOCUMENTO
\begin{document}

\maketitle
\thispagestyle{empty}

\begin{abstract}


\end{abstract}

\newpage
\tableofcontents

\newpage
\section*{Introducción}
\addcontentsline{toc}{section}{Introducción}

Este Proyecto busca implementar el cliente y el servidor de un chat multiusuario
funcional mediante la implementación de ciertos patrones de diseño que faciliten
su desarrollo. A su vez, se espera seguir el enfoque de gestión de proyectos
usando el modelo de la cascada, así como prácticas como las pruebas unitarias o
el código limpio, que no solo facilitan el desarrollo de proyectos, sino que
aseguran una mejor comprensión de cualquier proceso que lo implemente.

\section*{Problema}
\addcontentsline{toc}{section}{Problema}

Al ser un chat multiusuario, como su nombre lo indica, se espera que varias
personas se conecten al servidor de manera simultánea sin generar problemas
tanto para el cliente como para el servidor. En esencia, el servidor debe
permanecer siempre activo, ningún error debe terminar el programa, y en caso de
que alguno se detecte, el servidor debe ser capaz de manejarlo sin que afecte a
ninguno de los usuarios.

Para llevar esto a cabo se necesita del uso de programación concurrente, esto
para poder realizar diferentes tareas sin darle prioridad a ninguna en
particular.

También es necesaria una comunicación entre el cliente y el servidor. Los
usuarios no pueden comunicarse directamente entre sí; en su lugar, necesitan
enviar una solicitud al servidor para que este la procese y genere una
respuesta. Este proceso se aplica a cualquier acción que el usuario desee
realizar. Primero tiene que pasar por el servidor y, dependiendo de su validez,
se realizará.

\section*{Análisis}
\addcontentsline{toc}{section}{Análisis}

A grandes rasgos, este chat se puede dividir en el cliente y el servidor,
simplemente explicando las acciones que realizan y cómo interactúan entre sí a
través de los mensajes que se envían.

\begin{itemize}
    \newItem{Identificación del usuario}{Solicitud para conectarse e
        identificarse con un nombre de usuario.}{Si el nombre está disponible,
        se envía un mensaje al resto de los clientes informando del nuevo
        usuario, en caso contrario responderá que el nombre de usuario ya está
        en uso.}
    \newItem{Cambiar estado}{Cambiar el estado del usuario entre \texttt{active},
        \texttt{away} y \texttt{busy}.}{Si el nuevo estado es diferente al
        anterior, el servidor informa a los demás clientes del cambio.}
    \newItem{Lista de usuarios}{Solicita la lista de los usuarios en el
        chat.}{Regresa un objeto con los nombres y estados de los usuarios.}
    \newItem{Mensaje privado}{Mandar un mensaje privado a un usuario.}{Si el
        usuario existe el mensaje se envía, en caso contrario responderá que el
        usuario no existe.}
    \newItem{Mensaje público}{Mandar un mensaje público a todos los usuarios
        disponibles.}{No responde y envía el mensaje a todos los usuarios.}
    \newItem{Crear sala de chat}{Crear una nueva sala de chat con un nombre.}{Si
        el nombre está disponible, se creará una sala con solo el usuario en
        ella, en caso contrario responderá que la sala ya existe.}
    \newItem{Invitar usuarios}{Invitar a uno o más usuarios a una sala a la que
        pertenece el usuario.}{Si la sala y los usuarios invitados existen, el
        servidor no responde e invita a cada usuario, en caso contrario
        responderá que la sala (o el usuario) no existe.}
    \newItem{Unirse a una sala}{Unirse a una sala a la que el usuario ha sido
        invitado.}{Si la sala existe y el usuario fue invitado, el usuario se
        unirá a la sala, en caso contrario el servidor responderá que el cuarto
        no existe o que el usuario no fue invitado, según sea el caso.}
    \newItem{Usuarios en la sala}{Solicita un diccionario con los usuarios de la
        sala.}{Si la sala existe y el usuario pertenece a ella, el servidor
        regresará un diccionario con los nombres y estados de los usuarios. Si la
        sala existe pero el usuario no ha sido invitado o no se ha unido, el
        servidor responderá que el usuario no se ha unido a la sala y, si la sala
        no existe, el servidor transmitirá el mensaje.}
    \newItem{Mensaje a una sala}{Mandar un mensaje a una sala}{Si la sala existe
        y el usuario pertenece a ella, el servidor no responderá y envía el
        mensaje a todos los usuarios de la sala. Si la sala existe pero el
        usuario no ha sido invitado o no se ha unido, el servidor responderá que
        el usuario no se ha unido a la sala y, si la sala no existe, el servidor
        transmitirá el mensaje.}
    \newItem{Abandonar una sala}{Abandar una sala a la que pertenece el
        usuario.}{Si la sala existe y el usuario pertenece a ella, el servidor no
        responderá y el usuario abandonara la sala. Si la sala existe pero el
        usuario no ha sido invitado o no se ha unido, el servidor responderá que
        el usuario no se ha unido a la sala y si la sala no existe el servidor
        transmitirá el mensaje.}
    \newItem{Desconectarse}{Desconectar al usuario del chat y abandonar todas
        las salas a las que pertenece.}{No responde y el usuario se
        desconecta. Si este pertenece a alguna sala, las abandona.}
\end{itemize}


\section*{Diseño}
\addcontentsline{toc}{section}{Diseño}

\uml{Usuario}{
    \private{nombre: string}
    \private{estado: Estado}
    \private{color: Color}
}{
    \public{Usuario(nombre : string)}
    \public{getNombre(): string}
    \public{setEstado(estado : Estado)}
    \public{getEstado(): Estado}
    \public{getColor(): Color}
}

\uml{«enumeration»\\ Estado}{
    \noAccess{ACTIVE}
    \noAccess{AWAY}
    \noAccess{BUSY}
}{
    \noAccess{}
}

\uml{Color}{
    \private{colores : string[*]}
}{
    \public{random()}
}



\uml{Cliente}{
    \private{Usuario}
}{
    \public{hacerSolicitud()}
    \public{conectarse()}
    \public{cambiarEstado()}
    \public{listaUsuarios()}
    \public{mensajePrivado()}
    \public{mensajePublico()}
    \public{nuevaSala()}
    \public{invitarUsuarios()}
}

\uml{Servidor}{
    \private{Conexiones}
    \public
}{
    \public{responderSolicitud()}
    \public{usuarioValido()}
    \public{listaUsuarios()}
    \public{enviarMensaje()}
    \public{salaValida()}
    \public{crearSala()}

}

\uml{«enumeration»\\ Mensaje}{
    \noAccess
}{
    \public
}


\section*{Conclusión}
\addcontentsline{toc}{section}{Conclusión}



% FIN DEL DOCUMENTO
\end{document}
